\documentclass[12pt]{article}
\usepackage[margin=1in]{geometry}
\usepackage{amsmath,amssymb,mathtools}
\usepackage{enumitem}
\usepackage{bm}
\usepackage{hyperref}
\usepackage{fancyhdr}
\usepackage{draftwatermark}
\usepackage{xcolor}
\usepackage{array}
\usepackage{booktabs}

%------------- TITLE AND METADATA -------------

\newcommand{\CourseCode}{HE2001 }
\newcommand{\CourseName}{Microeconomics II}
\newcommand{\DocType}{Midterm Practice 2 (Questions)}

\title{
\CourseCode \CourseName \\[0.3em]
\large \DocType
}
\author{
  Academic Year 2025/2026, Semester 2 \\[0.25em]
  \textit{Quantitative Research Society @NTU}
}
\date{\today}

%----------------------------------------------

%------------- HEADER & FOOTER (for page 2 onward) -------------
\fancypagestyle{main}{
  \fancyhf{}
  \fancyhead[L]{\CourseCode - \DocType}
  \fancyhead[R]{\href{[https://qrsntu.org](https://qrsntu.org)}{QRS@NTU}}
  \fancyfoot[C]{\thepage}
  \renewcommand{\headrulewidth}{0.4pt}
  \renewcommand{\footrulewidth}{0pt}
}


%------------- SIMPLE FOOTER (page 1 only) -------------
\fancypagestyle{plainfirst}{
  \fancyhf{}
  \renewcommand{\headrulewidth}{0pt}
  \renewcommand{\footrulewidth}{0pt}
  \fancyfoot[C]{\scriptsize\textcolor{gray}{\textit{For academic reference only. This document is provided for educational purposes and does not constitute an official question paper, answer key, or any formally authorised academic material. Its content is not guaranteed to be complete or accurate.}}}
}
%------------------------------------------------------

%------------- WATERMARK -------------
\SetWatermarkText{Unofficial — QRS@NTU — qrsntu.org}
\SetWatermarkScale{1}
\SetWatermarkLightness{0.99}
%------------------------------------

\begin{document}

%------------- COVER PAGE -------------
\maketitle
\thispagestyle{plainfirst}
\pagestyle{main}
\bigskip

\section*{Instructions}

\begin{itemize}
\item This exam consists of 5 short questions (8 marks each) and 3 long questions (20 marks each).
\item Total marks: 100.
\item Show all working clearly. Answers without justification may receive zero marks.
\item The notation $(x_1, x_2)$ denotes a bundle of goods 1 and 2.
\item The notation $(p_1, p_2)$ denotes prices of goods 1 and 2.
\end{itemize}

\bigskip

\newpage
%----------------------------
\part{Short Questions}
%----------------------------

\section*{Question 1 (8 marks)}

Alice maximizes a strictly increasing utility function. You observe the following choices:
\begin{itemize}
\item At prices $p^1=(2,1)$, she buys $x^1=(4,2)$
\item At prices $p^2=(1,2)$, she buys $x^2=(2,4)$
\item At prices $p^3=(1,1)$, she buys $x^3=(3,3)$
\end{itemize}

\begin{enumerate}[label=(\alph*)]
\item (2 marks) Compute expenditure in each observation.
\item (3 marks) Determine all nontrivial direct revealed preference relations and strict direct revealed preference relations among $x^1,x^2,x^3$.
\item (3 marks) Does the data satisfy WARP? Does it satisfy GARP? Clearly justify using $R$, $P$, and (if needed) $R^*$.
\end{enumerate}

\bigskip

\section*{Question 2 (8 marks)}

A worker has utility
\[
U(c,\ell)=c^{1/2}\ell^{1/2},
\]
where $c$ is consumption and $\ell$ is leisure (hours). The worker has 16 hours per day, wage $w>0$, and non-labour income $y=40$.

\begin{enumerate}[label=(\alph*)]
\item (3 marks) Write the budget constraint and derive the optimal labour supply function $L^*(w)$ (hours worked), assuming an interior solution.
\item (2 marks) Compute optimal labour supply and consumption when $w=10$.
\item (3 marks) The wage rises to $w'=20$. Using the Slutsky method (compensating non-labour income so the original bundle remains affordable), compute the substitution effect and income effect on labour supply.
\end{enumerate}

\bigskip

\section*{Question 3 (8 marks)}

Consider a three-period economy $(t=0,1,2)$ with interest rate $r=0.25$.

\begin{enumerate}[label=(\alph*)]
\item (2 marks) Find the present value at $t=0$ of the cashflow stream $(0,125,156.25)$.
\item (3 marks) A project costs $210$ today and pays $(0,50,250)$ in periods $(0,1,2)$ respectively (so the $t=0$ component is the cost). Compute its NPV and state whether it should be accepted.
\item (3 marks) A perpetuity pays $30$ every period starting from $t=1$. If its market price is $130$, determine whether an arbitrage opportunity exists (assume short selling is allowed) and describe a valid arbitrage strategy.
\end{enumerate}

\bigskip

\section*{Question 4 (8 marks)}

There are four equally likely states. Consider two contingent consumption bundles (written in ascending order of states):
\[
x=(2,4,8,10),\qquad y=(1,5,7,11).
\]

\begin{enumerate}[label=(\alph*)]
\item (2 marks) Compute $\mathbb E[x]$ and $\mathbb E[y]$.
\item (3 marks) Does $x$ first-order stochastically dominate $y$? Does $y$ first-order stochastically dominate $x$? Justify using CDF comparisons at the relevant support points.
\item (3 marks) Determine whether $x$ second-order stochastically dominates $y$ (or vice versa). Use an appropriate discrete SSD test.
\end{enumerate}

\bigskip

\section*{Question 5 (8 marks)}

There are two states of nature, state 1 and state 2. Two assets are traded:
\begin{itemize}
\item Asset A pays $(4,0)$ and costs $1.00$
\item Asset B pays $(1,2)$ and costs $1.25$
\end{itemize}

\begin{enumerate}[label=(\alph*)]
\item (3 marks) Find the state price vector $(q_1,q_2)$. Also find the price and gross return of a risk-free asset paying $(1,1)$.
\item (2 marks) What is the no-arbitrage price of the contingent claim paying $(3,5)$?
\item (3 marks) A new asset pays $(2,1)$. What is its no-arbitrage price? If its market price is $0.90$, describe an arbitrage strategy.
\end{enumerate}

\newpage
%----------------------------
\part{Long Questions}
%----------------------------

\section*{Question 6 (20 marks)}

Consider a consumer with utility function
\[
U(x_1,x_2)=\ln x_1+2\ln x_2,
\]
with $x_1>0$ and $x_2>0$.

\begin{enumerate}[label=(\alph*)]
\item (4 marks) Compute $MU_1$, $MU_2$, and the marginal rate of substitution $MRS_{1,2}$. State the interior first-order condition for utility maximization at prices $(p_1,p_2)$.
\item (5 marks) Derive the Marshallian demand functions $x_1(p_1,p_2,m)$ and $x_2(p_1,p_2,m)$.
\item (4 marks) The consumer has endowment $\omega=(9,6)$ and faces prices $(p_1,p_2)=(2,3)$. Compute the value of the endowment, the optimal consumption bundle, and whether the consumer is a net buyer or seller of each good.
\item (3 marks) Verify Walras' Law at the bundle obtained in part (c).
\item (4 marks) Holding nominal income fixed at the value from part (c), suppose the price of good 1 rises from $p_1=2$ to $p_1'=3$ (with $p_2=3$ unchanged). Using the Slutsky decomposition for good 1, compute the substitution effect and income effect, and state whether good 1 is normal.
\end{enumerate}

\medskip

\newpage
\section*{Question 7 (20 marks)}

Consider a two-period intertemporal consumption problem with utility
\[
U(x_0,x_1)=\ln x_0+\ln x_1,
\]
where $x_0$ is period-0 consumption and $x_1$ is period-1 consumption. The consumer can borrow or save at interest rate $r$.

\begin{enumerate}[label=(\alph*)]
\item (3 marks) Write the intertemporal budget constraint in present-value form and in future-value form.
\item (5 marks) Let $m$ denote the present value of lifetime income. Derive the optimal demands $x_0^*(r,m)$ and $x_1^*(r,m)$.
\item (4 marks) Suppose the income stream is $(y_0,y_1)=(18,30)$ and $r=0.25$. Compute the present value of income, the optimal consumption bundle, and whether the consumer is a net borrower or net saver in period 0 (and by how much).
\item (4 marks) Verify that the optimal bundle from part (c) satisfies the budget constraint exactly.
\item (4 marks) Now suppose the interest rate rises to $r'=0.50$. Using Slutsky decomposition for period-0 consumption $x_0$, compute the Slutsky compensated present value $m^s$, the substitution effect, and the income effect. Is period-0 consumption normal?
\end{enumerate}

\bigskip

\newpage
\section*{Question 8 (20 marks)}

There are three states of nature $s=1,2,3$. Three traded assets have the following state payoffs and prices:
\[
\begin{array}{c|ccc|c}
\text{Asset} & s=1 & s=2 & s=3 & \text{Price} \\
\hline
1 & 2 & 0 & 0 & 0.40 \\
2 & 1 & 1 & 0 & 0.50 \\
3 & 1 & 1 & 1 & 0.90
\end{array}
\]

A consumer has expected utility
\[
U(x_1,x_2,x_3)=\frac12\ln x_1+\frac13\ln x_2+\frac16\ln x_3,
\]
and endowment $\omega=(4,2,1)$.

\begin{enumerate}[label=(\alph*)]
\item (5 marks) Find the state price vector $(q_1,q_2,q_3)$. Then compute the price and gross return of the risk-free payoff $(1,1,1)$.
\item (3 marks) Compute the value of the endowment and the no-arbitrage price of the contingent claim $d=(3,1,2)$.
\item (6 marks) Using the state-price budget constraint, solve for the consumer's optimal contingent consumption bundle $(x_1^*,x_2^*,x_3^*)$.
\item (3 marks) Find a portfolio $(\theta_1,\theta_2,\theta_3)$ of the traded assets that replicates the optimal bundle from part (c), and verify that its cost equals the state-price cost of the bundle.
\item (3 marks) Suppose the claim $d=(3,1,2)$ from part (b) is instead traded at price $1.80$. Is there an arbitrage opportunity? If yes, describe one and compute the initial arbitrage profit.
\end{enumerate}

\end{document}